\section*{Теоретическая справка}
\subsection*{Таблица первообразных элементарных функций}
Функция $F(x)$ называется первообразной для функции $f(x)$ на заданном промежутке, если для всех $x$ из этого промежутка выполняется равенство $F'(x)=f(x)$. Неопределённый интеграл $\int f(x)\,dx=F(x)+C$ --- множество всех первообразных функций $f(x)$.

\begin{center}
\renewcommand{\arraystretch}{1.35}
\begin{tabular}{c|c@{\qquad}c|c}
Функция $f(x)$ & Первообразная $F(x)$ & Функция $f(x)$ & Первообразная $F(x)$ \\ \hline
$c$ & $cx$ & $\sin x$ & $-\cos x$ \\
$x^n$, $n\ne-1$ & $\dfrac{x^{n+1}}{n+1}$ & $\cos x$ & $\sin x$ \\
$\sqrt{x}$ & $\dfrac23x\sqrt{x}$ & $\dfrac1{\cos^2x}$ & $\tg x$ \\
$\dfrac1{\sqrt{x}}$ & $2\sqrt{x}$ & $\dfrac1{\sin^2x}$ & $-\ctg x$ \\
$\dfrac1x$ & $\ln|x|$ & $e^x$ & $e^x$ \\
$\dfrac1{\sqrt{1-x^2}}$ & $\arcsin x$ & $a^x$ & $\dfrac{a^x}{\ln a}$ \\
$-\dfrac1{\sqrt{1-x^2}}$ & $\arccos x$ & $\dfrac1{1+x^2}$ & $\arctg x$ \\
$-\dfrac1{1+x^2}$ & $\arcctg x$ & &
\end{tabular}
\end{center}

\subsection*{Правила вычисления неопределённого интеграла}
\begin{enumerate}
  \item Интеграл от суммы или разности равен сумме или разности интегралов:
  \[\int(f(x)\pm g(x))\,dx=\int f(x)\,dx\pm\int g(x)\,dx=F(x)\pm G(x)+C.\]
  \item Постоянный множитель можно выносить за знак интеграла:
  \[\int a\cdot f(x)\,dx=a\int f(x)\,dx=aF(x)+C.\]
  \item Интегрирование сложной функции с линейным аргументом:
  \[\int f(kx+b)\,dx=\frac1kF(kx+b)+C.\]
\end{enumerate}

\subsection*{Метод замены переменной}
Если подынтегральное выражение не является табличным, часто помогает введение новой переменной. Например,
\[
\int\frac{dx}{x(1+\ln^2x)}.
\]
Положим $\ln x=t$, тогда $\frac1x\,dx=dt$ и
\[
\int\frac{dt}{1+t^2}=\arctg t+C=\arctg(\ln x)+C.
\]

\subsection*{Метод интегрирования по частям}
Из формулы производной произведения $(UV)'=U'V+UV'$ получаем в дифференциалах $d(UV)=V\,dU+U\,dV$, а значит:
\[
  \int U\,dV=U\cdot V-\int V\,dU.
\]
Мнемоническое правило: \textit{«УДиВился У Всех на ВиДУ»}.

\subsection*{Определённый интеграл и интегральная сумма}
Пусть функция $f(x)$ определена на отрезке $[a;b]$. Разобьём его на $n$ равных частей точками $a=x_0<x_1<\dots<x_n=b$ и выберем на каждом отрезке точку $c_i$. Интегральной суммой называется
\[
f(c_1)\Delta x+f(c_2)\Delta x+\dots+f(c_n)\Delta x.
\]
Геометрически это сумма площадей прямоугольников под графиком. Определённый интеграл --- предел интегральной суммы при $\Delta x\to0$:
\[
\int_a^b f(x)\,dx.
\]
Для вычисления служит формула Ньютона--Лейбница:
\[
\int_a^b f(x)\,dx=F(b)-F(a),
\]
где $F(x)$ --- любая первообразная функции $f(x)$.
